\section{Store a fraction as a scaled integer}
A ball can move by less than one pixel per simulation step. If the program stores
only the integer pixel coordinate, an increment of three quarters of a pixel has
no representation in that state. Rounding that increment to zero makes the ball
stationary; rounding it to one makes the ball too fast. Retain a more precise
position and choose a display cell when rendering it.

This chapter represents each coordinate as a signed 32-bit integer with eight
fractional bits. A raw integer \(r\) represents
\[
  x=\frac{r}{256},\qquad 256=2^8.
\]
One pixel is 256 raw units, and one raw unit is \(1/256\) pixel, or 0.00390625
pixel. The binary point is a convention used by the program. The register and
memory word contain an integer bit pattern; neither marks a binary-point position.

For this exercise the drawing is an eight-by-eight grid. Horizontal positions
increase rightward, and vertical positions increase downward. The ball is represented
by one selected cell, with its model position measured from the center
of cell zero to the center of cell seven on each axis. Each position is therefore
between zero and seven pixels, inclusive, or between zero and 1792 raw units.
Velocity is measured in pixels per simulation step. A raw velocity of 192 means
\(192/256=3/4\) pixel per step. This is a discrete simulation clock; a step does
not promise a particular duration in host seconds.

\Needspace{14\baselineskip}
\tagpdfsetup{table/header-rows={1}}
\begin{longtable}{rrl}
  \toprule
  Raw integer & Represented value & Role in the walkthrough     \\
  \midrule
  1           & \(1/256\)         & Smallest positive increment \\
  128         & \(1/2\)           & Half a pixel                \\
  192         & \(3/4\)           & Initial horizontal velocity \\
  256         & 1                 & One pixel                   \\
  320         & \(1+1/4\)         & Initial horizontal position \\
  448         & \(1+3/4\)         & Initial vertical position   \\
  1792        & 7                 & Upper coordinate bound      \\
  \bottomrule
\end{longtable}
To encode an exactly representable position, multiply its pixel value by 256.
For example, \(1.25\cdot256=320\). A value whose product is not an integer needs
an explicit input-quantization rule. The initial data in this exercise already
consist of raw integers, so no decimal-to-fixed conversion occurs during execution.

\section{Keep the whole raw word signed}
Eight fractional bits do not make the low byte a separate signed fraction.
The whole 32-bit raw integer uses two's-complement signed interpretation. Raw
\(-96\), stored as hexadecimal \texttt{ffffffa0}, represents \(-3/8\).
Its low byte is hexadecimal \texttt{a0}, or 160. Reading that byte as an independent
positive fraction would lose the meaning of the higher bits.

An arithmetic right shift by eight rounds a negative raw value downward when
interpreted as division by 256. For raw \(-96\), that integer part is \(-1\),
and the low eight bits supply the remaining \(160/256\):
\[
  -1+\frac{160}{256}=-\frac{3}{8}.
\]
Signed integer division truncated toward zero would instead produce zero as its
quotient. These are different rounding rules. The renderer below receives only
nonnegative positions, but velocities can be negative. Keep that distinction when
choosing a shift or interpreting a register.

With a signed 32-bit raw word, the representable range is
\[
  -2^{23}\leq x\leq2^{23}-\frac1{256}.
\]
Descriptions using the letter Q differ in whether the sign bit is counted among
the integer bits. Stating the total width and the eight fractional bits avoids
that ambiguity. The ball uses only a small bounded part of this format's range.
